\section{Introduction}

A smooth projective variety is universally \(CH_0\)-trivial if the degree map
\[
  \deg:\CH_0(Y_L)\longrightarrow\Z
\]
is an isomorphism for every field extension \(L/\C\).  This condition is a
central consequence of stable rationality, and its failure is a standard
stable-rationality obstruction.  Its validity, however, need not provide a
rational parametrization.  We ask how far that gap can persist after one
explicit stabilization.

The irrationality half of that question has a uniform answer.  The small even
quantum connection of every smooth cubic threefold has primitive-sixth framed
formal monodromy, and Iritani's quantum blowup formula shows that its
multiplicity survives one stabilization.

\begin{theorem}[One-step irrationality]
\label{thm:every-cubic}
For every smooth complex cubic threefold \(X\), the fourfold
\(X\times\PP^1\) is irrational.
\end{theorem}

The same invariant also gives one-step irrationality for every smooth prime
Fano threefold of genus eight; see Corollary~\ref{cor:v14-one-step}.

Yang--Yu--Zhu construct a two-dimensional family of smooth cubic threefolds
admitting unirational parametrizations of degrees two and three
\cite[Theorems 1.1 and 3.3]{YYZ}.  Unirational parametrizations of coprime
degrees give a decomposition of the diagonal, so every member of that family
is universally \(CH_0\)-trivial, and the same parametrizations persist on
\(X\times\PP^m\) \cite[Corollary~3.5]{YYZ}.  Together with
Theorem~\ref{thm:every-cubic}, their family already answers the question
above on a two-dimensional locus in moduli.

Our second result exhibits the separation on a family whose universal
\(CH_0\)-triviality is proved from the intermediate Jacobian instead.  Here
\(X\) is a smooth cubic threefold in the nonstandard \(A_5\)-invariant
pencil, and the primitive minimal class of its intermediate Jacobian is
algebraic.

\begin{theorem}[Separation on the \(A_5\)-pencil]
\label{thm:separation-family}
Let \(W_5\) be the five-dimensional irreducible \(A_5\)-module, and let
\(B^\circ\) be the smooth locus inside
\[
 \PP\bigl((\operatorname{Sym}^3W_5^*)^{A_5}\bigr).
\]
For the resulting family \(\mathcal X\to B^\circ\) and every
 \(b\in B^\circ(\C)\),
\[
 \mathcal X_b\text{ is universally }CH_0\text{-trivial},
 \qquad
 \mathcal X_b\times\PP^1\text{ is universally }CH_0\text{-trivial and irrational}.
\]
The family is non-isotrivial in coarse moduli.
\end{theorem}

The invariant cubic space in Theorem~\ref{thm:separation-family} is
two-dimensional.  Indeed, the character of \(W_5\) is
\((5,1,-1,0,0)\) on the classes \(1,2A,3A,5A,5B\).  The symmetric-cube
character formula gives
\[
 \dim \bigl(\operatorname{Sym}^3W_5^*\bigr)^{A_5}=2.
\]
Thus its projectivization is a pencil with connected smooth locus
\(B^\circ\).  Hartlieb denotes it by \(M_{H_1}\); it is not the component
associated with \(\mathbf1\oplus W_4\) \cite[Lemma~5.5]{Hartlieb}.  Since the
family is non-isotrivial, its image in coarse moduli is a curve, one
dimension below the locus of \cite{YYZ}.

\begin{remark}
\label{rem:fermat-in-pencil}
The pencil contains the Fermat cubic threefold.  Let \(A_4\subset A_5\) be a
point stabilizer of the degree-five action and \(\chi\) a nontrivial cubic
character of \(A_4\).  Frobenius reciprocity gives
\[
 \bigl\langle\operatorname{Ind}_{A_4}^{A_5}\chi,\;W_5\bigr\rangle
 =\bigl\langle\chi,\;\operatorname{Res}_{A_4}W_5\bigr\rangle=1,
\]
and both sides of the induction have dimension five, so
\(\operatorname{Ind}_{A_4}^{A_5}\chi\cong W_5\).  Induction from a character
realizes \(W_5\) by monomial matrices whose nonzero entries are cube roots of
unity, and such a matrix fixes \(y_1^3+\cdots+y_5^3\).  For that member of the
pencil, universal \(CH_0\)-triviality is already known three times over: from
Voisin's minimal-class examples \cite[Theorem~4.5]{Voisin}, from
Colliot-Th\'el\`ene's almost-diagonal theorem \cite{CT}, and from the
degree-three unirational parametrization of \cite[Remark~3.6]{YYZ}.
\end{remark}

We do not claim that \(X\) is stably irrational.  From the second
stabilization onward, weak-factorization centers can themselves have cubic
dimension, and the low-dimensional exclusion used here no longer applies.

\subsection*{How the two proofs meet}

The cycle argument starts from six elliptic quotients of the Fano surface.
Their Rosati Gram matrix is \(6I-J\), and the geometric kernel lies in the
exotic two-point complement of \(\PP^1(\F_2)\subset\PP^1(\F_4)\).  Its local
graph slopes are finite etale at two and three.  The graph equations then
make the descended divisor lattices rank-one generated.  Since rank-one
coefficient divisors have square zero, \(\Theta^4/4!\) is an ordinary
integral divisor product.

The quantum argument counts primitive sixth-root eigenvalues of framed formal
monodromy on the numerical small even quantum connection.  Iritani's
operation formulas make this count additive through weak factorization;
divisor tags repair a potentially noninjective center specialization.
Starting from Cai's displayed small-even matrix, the internal calculation gives
the nonzero count, while every center of dimension at most
two contributes zero.  The special \(A_5\)-geometry enters only in the cycle
argument; the quantum theorem applies to every smooth cubic threefold.

\subsection*{Relation with earlier rationality results}

Clemens and Griffiths proved that every smooth complex cubic threefold is
irrational by means of its intermediate Jacobian \cite{CG}.  That obstruction
does not by itself control a product with projective space, and stable
rationality of smooth cubic threefolds remains open.  Cai recently recovered
irrationality from the primitive-sixth-root formal monodromy of the cubic
quantum connection \cite{Cai}.  Iritani and Iritani--Koto prove the quantum
\(D\)-module decompositions for blowups and projective bundles
\cite{IritaniBlowup,IritaniKoto}; Katzarkov--Kontsevich--Pantev--Yu identify
the nef-canonical normal form used for low-dimensional centers
\cite{KKPYY}.  Theorem~\ref{thm:every-cubic} combines these inputs at the next
dimension.

On the cycle side, Voisin reduces universal \(CH_0\)-triviality of a cubic
threefold to algebraicity of the primitive minimal class of its intermediate
Jacobian \cite{Voisin}.  Beckmann--de Gaay Fortman place such minimal classes
in a general integral-Fourier framework \cite{BdGF}.  For the very general
cubic threefold the minimal class is not algebraic: every curve class on the
intermediate Jacobian is an even multiple of it, so a very general cubic
threefold admits no decomposition of the diagonal
\cite[Theorem~1.3 and Corollary~1.4]{EdGFS}.  What Voisin's criterion leaves is
the countable union of moduli subvarieties on which the class is algebraic.
Voisin shows that this union has components of codimension at most three and
exhibits explicit members \cite[Theorem~4.5 and Lemma~4.6]{Voisin}, and
Colliot-Th\'el\`ene proves universal \(CH_0\)-triviality for every smooth
cubic hypersurface whose equation is a sum of forms in separated groups of at
most three variables \cite{CT}.  The coprime-degree family of \cite{YYZ} lies
in the same union, by a decomposition of the diagonal rather than by a period
computation.  Theorem~\ref{thm:six-axis-divided-powers} does not enlarge the
union.  It proves algebraicity uniformly along one explicit pencil, from the
six-axis lattice rather than from an isogeny to a curve Jacobian or from a
unirational parametrization, and that pencil meets the earlier examples in
the Fermat cubic threefold (Remark~\ref{rem:fermat-in-pencil}).  The six elliptic
quotients supplied by Roulleau and the special period locus isolated by
Hartlieb make the class explicit: it belongs to the
ordinary integral divisor-product lattice for every smooth member of the
nonstandard \(A_5\)-pencil.

\subsection*{Proof map}

Section~\ref{sec:envelope} computes the six-axis polarization and identifies
the exotic cubic gluing.  Section~\ref{sec:minimal-class} proves
the local divisor-product theorem and applies it to the primitive minimal
class.  Section~\ref{sec:one-step} proves Theorem~\ref{thm:every-cubic} from
formal monodromy and weak factorization.  The final section combines the two
detectors and records the exact scope boundary.
